\section{Atomic descent and finite-pool inexactness}
\label{sec:descent}

This section converts the distinctions of \Cref{sec:setting} into energy inequalities. The standard
atomic residual estimate is derived first, since its orthogonality hypothesis is weakened in the
finite-tolerance algorithm.

\begin{lemma}[Atomic residual bound]
\label{lem:atomic-residual}
Under \Cref{ass:atomic} and the exact orthogonality \cref{eq:orthogonality},
\begin{equation}
  \norm{\E'(G_m)}_{\D^*}\geq \frac{\delta_m}{B}.
\label{eq:atomic-residual}
\end{equation}
\end{lemma}

\begin{proof}[Proof sketch]
Convexity gives the atomic residual estimate, the unit-step quadratic bound gives the reciprocal recurrence, and strong convexity on the active space gives the KKT estimate.
\end{proof}

\begin{definition}[\(p\)-growth density]
\label{def:p-growth}
Let \(p\geq2\), write \(z=(s,\xi)\in\R\times\R^d\), and omit the scalar component \(s\) for a
pure-gradient energy. A twice differentiable convex density \(\Psi(x,z)\) has \(p\)-growth on the
relevant range if, for almost every \(x\),
\begin{align*}
 c_0\abs z^p-c_1&\leq \Psi(x,z)\leq c_2(1+\abs z^p),\\
 \norm{D_z^2\Psi(x,z)}&\leq c_3(1+\abs z)^{p-2},
\end{align*}
with positive constants independent of \(x\) and \(z\) in that range. The first line controls bounded
energy sublevels, whereas the Hessian bound is the part used in the directional estimate below.
\end{definition}

\begin{assumption}[Directional quadratic upper bound]
\label[assumption]{ass:directional}
There is \(L>0\) such that, on the trajectory sublevel set, for every \(g\in\D\) and
\(\abs\lambda\leq1\),
\begin{equation}
 \E(u+\lambda g)\leq\E(u)+\lambda\ip{\E'(u)}g+\frac L2\lambda^2.
\label{eq:directional-smoothness}
\end{equation}
\end{assumption}

For $a\geq0$, define the maximal decrease supplied by the directional quadratic
bound by
\[
 H_L(a):=\max_{0\leq\lambda\leq1}
       \left(\lambda a-\frac L2\lambda^2\right)
 =\begin{cases}a^2/(2L),&0\leq a\leq L,\\
                 a-L/2,&a>L.
   \end{cases}
\]

The assumption is directional; it does not require a globally Lipschitz derivative on the full Banach
space. For the standard \(p\)-energies, bounded \(W^{1,p}\) sublevels and an \(X\)-normalized dictionary
already give the needed estimate by H\"older's inequality; an \(L^\infty\) atom bound is only a
convenient stronger condition for more general densities. Indeed, for
\(F(\lambda)=\E(u+\lambda g)\), the diffusion part of \(F''\) is controlled by
\[
 C_p\int_\Omega(\abs{\nabla u}+\abs\lambda\abs{\nabla g}+\eps)^{p-2}
 \abs{\nabla g}^2\dd x,
\]
and H\"older's inequality gives
\(F''(\lambda)\leq C_p(\norm{\nabla u}_{L^p}+\norm{\nabla g}_{L^p}+\eps|\Omega|^{1/p})^{p-2}
\norm{\nabla g}_{L^p}^2\) for \(\abs\lambda\leq1\). A reaction term is handled similarly. Removing
almost-null projected atoms remains important for the numerical parameterization and for uniform
source/entropy constants, but it is not necessary for this abstract \(L^p\) estimate once the atoms are
already normalized in the energy space.

\begin{theorem}[Generic energy baseline]
\label{thm:generic}
Suppose \Cref{ass:atomic,ass:directional} hold and \(t_m\geq t>0\). Then the exact continuous CGA
satisfies
\begin{equation}
  \delta_m\leq \frac{C(B,L,t,\delta_0)}{m+1}.
\label{eq:generic-rate}
\end{equation}
\end{theorem}

\begin{proof}[Proof sketch]
Convexity gives the atomic residual estimate, the unit-step quadratic bound gives the reciprocal recurrence, and strong convexity on the active space gives the KKT estimate.
\end{proof}

The standard $p$-energy setting used below is
\[
 \E_p(u)=\int_\Omega \frac1p|\nabla u|^p\,dx-\ell(u),\qquad
 \E_{p,\varepsilon}(u)=\int_\Omega
 \frac{(\varepsilon^2+|\nabla u|^2)^{p/2}-\varepsilon^p}{p}\,dx-\ell(u),
\]
with the zero-mean space in the pure and regularized Neumann cases. With a reaction
term we use
\[
 \E_{p,R}(u)=\int_\Omega\frac{|\nabla u|^p+|u|^p}{p}\,dx-\ell(u)
\]
on $W^{1,p}(\Omega)$. The complete assumptions and spaces are specified in
Section~\ref{sec:pgeometry}.

\subsection{Inexact selection and correction}

Write \(\widehat\delta_m=\E(\widehat G_m)-\E(u^*)\). In place of exact orthogonality, suppose that
\begin{equation}
  \abs{\ip{\E'(\widehat G_m)}{\widehat G_m}}\leq\eta_m^{\rm stat}.
  \label{eq:orthogonality-defect}
\end{equation}
The decomposition contains one multiplicative resolution factor and three additive error
contributions: \(t_m^{\rm pool}\tau_m\) records weak selection and continuous-dictionary coverage,
\(\xi_m\) records score evaluation, \(\eta_m^{\rm stat}\) is the radial active-space stationarity
defect of \cref{eq:orthogonality-defect} alone, and \(\varepsilon_m^{\rm opt}\) records the energy
accuracy of the correction. The full active-space dual Karush--Kuhn--Tucker (KKT) residual is
denoted separately by \(r_m^{\mathrm{KKT}}\) below. In the terminology of the greedy-with-errors
framework of \cite[Section~3]{DereventsovTemlyakov2022}, \(t_m^{\rm pool}\tau_m\) is the effective
weakness parameter and \(\varepsilon_m^{\rm opt}\) is an energy-evaluation/correction error. The
quantities \(\xi_m\) and \(\eta_m^{\rm stat}\) separate the finite-score and biorthogonality
contributions that the abstract framework combines into a single stepwise inexactness allowance.

\begin{theorem}[Finite-pool one-step descent]
\label{thm:inexact}
Assume \Cref{ass:atomic,ass:directional}, \cref{eq:score-error,eq:pool-coverage}, and
\cref{eq:orthogonality-defect} on the numerical trajectory. Define
\[
 a_m=\left[
 t_m^{\rm pool}\tau_m\frac{\widehat\delta_m-\eta_m^{\rm stat}}{B}
 -(1+t_m^{\rm pool})\xi_m
 \right]_+.
\]
Here \([x]_+=\max\{x,0\}\).
Then
\begin{equation}
 \widehat\delta_{m+1}\leq\widehat\delta_m-H_L(a_m)+\varepsilon_m^{\rm opt}.
 \label{eq:inexact-descent}
\end{equation}
If, for a fixed absorption constant \(\vartheta\in(0,1)\),
\begin{equation}
 \eta_m^{\rm stat}+\frac{B(1+t_m^{\rm pool})}{t_m^{\rm pool}\tau_m}\xi_m
 \leq\vartheta\widehat\delta_m,
 \label{eq:relative-defect}
\end{equation}
the effective weakness \(\inf_{m\ge m_0} t_m^{\rm pool}\tau_m\) is positive, and there
exist a finite index \(m_0\) and \(0\leq\zeta<1\) such that
\[
  \varepsilon_m^{\rm opt}\leq\zeta\,H_L(a_m),\qquad m\ge m_0,
\]
then \(\widehat\delta_m=O((m-m_0+1)^{-1})\) for the tail. The finite prefix
\(m<m_0\) is not included in this rate statement.
\end{theorem}

\begin{proof}[Proof sketch]
Convexity gives the atomic residual estimate, the unit-step quadratic bound gives the reciprocal recurrence, and strong convexity on the active space gives the KKT estimate.
\end{proof}

\begin{remark}[Error saturation]
\Cref{eq:inexact-descent} shows how error saturation arises. A pool whose
\(\tau_m\) decreases, a score error \(\xi_m\) comparable with the best resolvable score, or a correction
error \(\varepsilon_m^{\rm opt}\) comparable with the predicted decrease can each stop the recursion. The equation
does not identify which defect dominates from a plotted error curve alone.
\end{remark}

The optimization error can be related to a KKT residual only after a norm and a strong-convexity
constant have been fixed. We say that $\E|_\Phi$ is $\mu_\Phi$-strongly convex in the norm
$\norm\cdot_\Phi$ if, for all $u,v\in\Phi$,
\[
 \E(v)\geq \E(u)+\langle\E'(u),v-u\rangle
       +\frac{\mu_\Phi}{2}\norm{v-u}_\Phi^2 .
\]
Here $\mu_\Phi$ is a constant on the finite-dimensional active space $\Phi$; it
cannot be extrapolated to a strong-convexity constant on the full space without an
additional argument.

\begin{lemma}[KKT residual to correction error]
\label{lem:kkt}
Let \(\Phi\) be finite dimensional and let \(\E|_\Phi\) be \(\mu_\Phi\)-strongly convex in a norm
\(\norm\cdot_\Phi\). If \(G_\Phi^{\sharp}\) is the exact minimizer and
\[
 r^{\mathrm{KKT}}=\norm{\E'(\widehat G)|_\Phi}_{\Phi^*},
\]
then
\begin{equation}
 0\leq\E(\widehat G)-\E(G_\Phi^{\sharp})
 \leq\frac{(r^{\mathrm{KKT}})^2}{2\mu_\Phi}.
\label{eq:kkt-energy}
\end{equation}
The same residual also satisfies
\begin{equation}
 \abs{\ip{\E'(\widehat G)}{\widehat G}}
 \leq r^{\mathrm{KKT}}\norm{\widehat G}_\Phi.
\label{eq:kkt-orth}
\end{equation}
\end{lemma}

\begin{proof}[Proof sketch]
Convexity gives the atomic residual estimate, the unit-step quadratic bound gives the reciprocal recurrence, and strong convexity on the active space gives the KKT estimate.
\end{proof}
